\section{随机变量是样本空间上的函数}
\label{sec:05-Probability/03-Random-Variables/01}

一次线上会话留下的原始记录是一长串东西：点击序列、每次停留时长、设备型号、请求时间戳。风控同事想知道这次会话可疑不可疑，产品同事想知道用户翻了几个页面，成本同事只关心这次会话烧了多少算力。三个人盯着同一批记录，各自读出一个数。

上一单元的语言在这里开始吃力。要问``页面数超过 5''的概率，得先在样本空间里圈出满足条件的那些会话；换成超过 6，再圈一遍。更麻烦的是，三位同事各自关心的那个数彼此有没有牵连，用集合语言几乎说不出来。

缺的零件很朴素：先把每条原始记录换算成一个数，再让概率跟着这个数走。

抛三枚硬币，八个等可能序列 HHH、HHT、HTH、THH、HTT、THT、TTH、TTT。问``有几个正面''，问的其实不是哪个序列出现了，而是一条从序列到数字的规则：HHH 记 3，HHT 记 2，TTT 记 0。换个问题就换条规则——``第一次是不是正面''给出 0 或 1，``最长连续正面有多长''又是另一条。同一批序列上，几条规则可以并存。

\begin{definition}
随机变量 \(X\) 是从样本空间 \(\Omega\) 到实数集 \(\R\) 的（可测）函数：
\[
X:\Omega\to\R.
\]
实验产生结果 \(\omega\)，函数吐出数值 \(X(\omega)\)。
\end{definition}

名字有点误导。\(X\) 一点也不变，变的是喂进去的 \(\omega\)：抛之前你不知道哪个序列会实现，抛完之后 \(X(\omega)\) 没有半点余地。所以``随机变量''其实是``定义在随机结果上的确定函数''，只是这个名字已经叫开了。做数据的人对此很熟——一列特征就是样本空间上的一个函数，特征工程做的事情，无非是在同一个 \(\Omega\) 上再挂几个函数。

\begin{center}
\begin{tikzpicture}[x=1cm,y=1cm,font=\footnotesize,>=stealth]
  \draw[black!45,fill=blue!5] (6.7,2.4) ellipse (6.7 and 1.05);
  \draw[black!35,fill=blue!16,rounded corners=6pt] (0.35,2.08) rectangle (10.55,2.72);
  \node at (1.0,2.4) {TTT};
  \node at (3.55,2.4) {HTT};
  \node at (4.8,2.4) {THT};
  \node at (6.05,2.4) {TTH};
  \node at (7.35,2.4) {HHT};
  \node at (8.6,2.4) {HTH};
  \node at (9.85,2.4) {THH};
  \node at (12.4,2.4) {HHH};
  \node[black!70] at (6.7,3.72) {\(\Omega\)：八个等可能序列，阴影部分是 \(\set{X\le 2}\) 的原像};
  \draw[->,black!45] (1.0,2.02) -- (1.0,0.46);
  \draw[->,black!45] (3.55,2.02) -- (4.8,0.46);
  \draw[->,black!45] (4.8,2.02) -- (4.8,0.46);
  \draw[->,black!45] (6.05,2.02) -- (4.8,0.46);
  \draw[->,black!45] (7.35,2.02) -- (8.6,0.46);
  \draw[->,black!45] (8.6,2.02) -- (8.6,0.46);
  \draw[->,black!45] (9.85,2.02) -- (8.6,0.46);
  \draw[->,black!45] (12.4,2.02) -- (12.4,0.46);
  \draw[blue!30,line width=2.4pt] (-0.1,0) -- (8.6,0);
  \draw[->] (-0.2,0) -- (13.5,0) node[right] {\(\R\)};
  \draw (1.0,-0.12) -- (1.0,0.12);
  \draw (4.8,-0.12) -- (4.8,0.12);
  \draw (8.6,-0.12) -- (8.6,0.12);
  \draw (12.4,-0.12) -- (12.4,0.12);
  \node[below] at (1.0,-0.16) {\(0\)};
  \node[below] at (4.8,-0.16) {\(1\)};
  \node[below] at (8.6,-0.16) {\(2\)};
  \node[below] at (12.4,-0.16) {\(3\)};
  \node[below,black!65] at (1.0,-0.62) {\(1/8\)};
  \node[below,black!65] at (4.8,-0.62) {\(3/8\)};
  \node[below,black!65] at (8.6,-0.62) {\(3/8\)};
  \node[below,black!65] at (12.4,-0.62) {\(1/8\)};
\end{tikzpicture}

\emph{\(X\) 把八个序列压到四个数上。数轴上加粗的那一段是集合 \((-\infty,2]\)；顺着箭头倒着走，它在 \(\Omega\) 里对应阴影中的七个序列，概率 \(7/8\)。底下一行分数是这次压缩留下的全部东西——\(X\) 的分布。}
\end{center}

\subsection{事件是数轴上集合的原像}

写 \(P(X\le 2)\) 时，括号里的 \(X\le 2\) 不是一次数值比较，它是
\[
\set{\omega\in\Omega:X(\omega)\le 2}
\]
的简写。一般地，对数轴上的（可测）集合 \(B\)，
\[
P(X\in B)=P\bigl(\set{\omega:X(\omega)\in B}\bigr).
\]
整个机制就这一行。概率测度还是原来那个 \(P\)，仍然长在 \(\Omega\) 上；随机变量只干一件事，把数轴上的集合 \(B\) 拉回样本空间，交给 \(P\) 去量。上图里那条加粗线段和那片阴影，是同一个事件的两种写法。

拉回操作在数轴上诱导出一份新的概率分配，称为 \(X\) 的分布。它才是后面所有计算真正打交道的对象：要回答关于正面数的任何问题，记住 \(1/8,3/8,3/8,1/8\) 这四个数就够了，八个序列可以退到幕后。取值不是整数也一样，令 \(W=X/3\) 表示正面比例，问 \(P(W>0.5)\)，先把它翻译成 \(X>1.5\)，再翻译成 \(X\in\set{2,3}\)，概率 \(4/8\)。翻译的每一步都在数轴上完成，样本空间只在最后一步被调用一次。

压缩当然有代价。\(X(\text{HHT})=X(\text{HTH})=X(\text{THH})=2\)，三个序列合成一团，正面出现在哪一位这件事被抹掉了。分布记得这一团总共占 \(3/8\)，不记得团内怎么分。所以``正面数为 2 时第一枚是不是正面''这个问题，分布回答不了——不是算不出来，是信息已经不在里面了。

\subsection{分布相同，两个变量可以完全不同}

掷一枚公平骰子，令 \(X\) 为点数，\(Y=7-X\)。两者的边缘分布一模一样，都是 \(1,\ldots,6\) 上的均匀分布。但它们不是同一个函数，也谈不上独立：\(X=3\) 一出现，\(Y=4\) 就是板上钉钉。

把另一枚独立骰子的点数记作 \(Z\)，边缘同样均匀。对比联合概率：
\[
P(X=3,Z=4)=\frac1{36},\qquad P(X=3,Y=4)=P(X=3)=\frac16.
\]
边缘一致，联合差了六倍。边缘分布只说每个变量各自取值的脾气，对它们怎样一起动弹一个字也没说，而``一起动弹''恰恰由它们同为 \(\Omega\) 上的函数这一点决定——不需要额外把它们``连接''起来，底层的 \(\omega\) 早就把取值绑在一起了。\hyperref[ch:041]{``联合分布：单个分布之外，变量怎样一起变化''}整章都在处理这层结构。

这条区分在模型评估里天天出现。两个模型在测试集上的误差分布几乎重合，不代表它们错在同一批样本上。集成有没有收益，看的是误差的联合行为；只看两条边缘的直方图，什么也判断不出来。

\subsection{指示变量把``发生了吗''换成算术}

对任意事件 \(A\)，最简单的随机变量是
\[
\ind_A(\omega)=\begin{cases}1,&\omega\in A,\\0,&\omega\notin A.\end{cases}
\]
它的用处在于当桥。事件的交对应指示变量的乘积（\(\ind_{A\cap B}=\ind_A\ind_B\)），事件的概率对应指示变量的期望（\(P(A)=\E[\ind_A]\)）。集合运算与算术运算这两套语言，从此可以互相翻译。

检查十封邮件，令 \(A_i\) 为``第 \(i\) 封是垃圾邮件''，垃圾邮件总数就是 \(X=\sum_{i=1}^{10}\ind_{A_i}\)。这个写法完全不要求各封邮件独立——求和是在单个 \(\omega\) 内部做的算术，把一串 0 和 1 加起来，与概率测度怎样分解无关。独立性要到你想把联合概率拆成乘积时才登场。\hyperref[sec:05-Probability/04-Expectation-and-Moments/01]{``期望、线性性与指示变量''}会兑现这笔账：期望的线性性让 \(\E[X]=\sum_i P(A_i)\) 无条件成立，哪怕垃圾邮件成堆出现。

\subsection{选了哪些函数，就决定了能问哪些问题}

随机变量划出一条分界线。线之前你在样本空间里摆弄具体的 \(\omega\)，线之后你只用数值和分布说话。真实的数据工作正是这么走的：面对一批病历、一串交易、一列传感器读数，你提取出年龄、金额、温度，此后所有推断都架在这些数值之上。

这条线也是风险所在。病历里没记吸烟史，那么无论逻辑回归拟合得多漂亮，它对肺癌风险的预测都会系统性偏低——关键信号从一开始就没被编码成任何一个函数。概率论负责从你看了的东西里得出结论，不负责提醒你漏看了什么。所以``选哪些随机变量''是建模里最早也最不可回退的一步。

\begin{remark}
严格地说 \(X\) 还必须可测：对每个实数 \(x\)，集合 \(\set{\omega:X(\omega)\le x}\) 得属于事件族 \(\mathcal F\)，否则 \(P(X\le x)\) 无从谈起。有限样本空间和常见的连续模型里，分段连续、单调、只有可数个间断点的函数全都自动过关，日常不必逐个检查。这条要求挡住的是这种事：借一个``函数''把不可测的怪集合偷渡回来，然后问出一个根本没有答案的概率问题。完整交代见\hyperref[ch:101]{``测度、积分与条件期望：为现代概率回填地基''}。
\end{remark}

分布现在是数轴上的一份概率分配，但还没有具体写法。取值可数时最省事：给每个可能出现的数标上一份质量，加起来正好是 1。下一节从这里入手。
